% ======================================================================
% APPENDIX A: MATHEMATICAL PROOFS & THEORETICAL DERIVATIONS
% ======================================================================
\chapter{Mathematical Proofs \& Theoretical Derivations}
\label{app:proofs}

\section{Proof of Bellman Optimality via Banach Fixed Point Theorem}
Let $\mathcal{B}(\mathcal{S} \times \mathcal{A})$ denote the Banach space of all bounded real-valued functions over state-action space $\mathcal{S} \times \mathcal{A}$, equipped with the supremum norm (infinity norm):
\begin{equation}
    \|Q\|_\infty = \sup_{s \in \mathcal{S}, a \in \mathcal{A}} |Q(s, a)|
\end{equation}
We define the Bellman Optimality Operator $\mathcal{T}^*: \mathcal{B}(\mathcal{S} \times \mathcal{A}) \to \mathcal{B}(\mathcal{S} \times \mathcal{A})$ by:
\begin{equation}
    (\mathcal{T}^* Q)(s, a) = \mathcal{R}(s, a) + \gamma \sum_{s' \in \mathcal{S}} \mathcal{P}(s' \mid s, a) \max_{a' \in \mathcal{A}} Q(s', a')
\end{equation}

\begin{theorem}[Contraction Mapping of Bellman Optimality Operator]
If the discount factor satisfies $\gamma \in [0, 1)$, then $\mathcal{T}^*$ is a $\gamma$-contraction mapping on $(\mathcal{B}(\mathcal{S} \times \mathcal{A}), \|\cdot\|_\infty)$.
\end{theorem}

\begin{proof}
Let $Q_1, Q_2 \in \mathcal{B}(\mathcal{S} \times \mathcal{A})$. For any state $s \in \mathcal{S}$ and action $a \in \mathcal{A}$:
\begin{align}
    |(\mathcal{T}^* Q_1)(s, a) - (\mathcal{T}^* Q_2)(s, a)| &= \left| \gamma \sum_{s'} \mathcal{P}(s' \mid s, a) \left( \max_{a'} Q_1(s', a') - \max_{a''} Q_2(s', a'') \right) \right| \nonumber \\
    &\le \gamma \sum_{s'} \mathcal{P}(s' \mid s, a) \left| \max_{a'} Q_1(s', a') - \max_{a''} Q_2(s', a'') \right|
\end{align}
Using the algebraic inequality $|\max_x f(x) - \max_x g(x)| \le \max_x |f(x) - g(x)|$:
\begin{align}
    |(\mathcal{T}^* Q_1)(s, a) - (\mathcal{T}^* Q_2)(s, a)| &\le \gamma \sum_{s'} \mathcal{P}(s' \mid s, a) \max_{a'} |Q_1(s', a') - Q_2(s', a')| \nonumber \\
    &\le \gamma \sum_{s'} \mathcal{P}(s' \mid s, a) \|Q_1 - Q_2\|_\infty \nonumber \\
    &= \gamma \|Q_1 - Q_2\|_\infty \sum_{s'} \mathcal{P}(s' \mid s, a) \nonumber \\
    &= \gamma \|Q_1 - Q_2\|_\infty
\end{align}
Taking the supremum over all $s \in \mathcal{S}$ and $a \in \mathcal{A}$:
\begin{equation}
    \|\mathcal{T}^* Q_1 - \mathcal{T}^* Q_2\|_\infty \le \gamma \|Q_1 - Q_2\|_\infty
\end{equation}
Since $\gamma < 1$, $\mathcal{T}^*$ is a strict contraction. By the Banach Fixed Point Theorem, there exists a unique fixed point $Q^* \in \mathcal{B}(\mathcal{S} \times \mathcal{A})$ such that $\mathcal{T}^* Q^* = Q^*$, and value iteration converges geometrically:
\begin{equation}
    \|Q_{k+1} - Q^*\|_\infty \le \gamma^k \|Q_0 - Q^*\|_\infty
\end{equation}
which proves asymptotic convergence to the optimal routing value function.
\end{proof}

\section{Derivation of Double Q-Learning Bias Elimination}
Consider a state $s$ with $M$ candidate actions whose true values are equal: $Q^*(s, a_i) = V^*$. Suppose our value estimates are subject to zero-mean random estimation noise $\epsilon_i$:
\begin{equation}
    \hat{Q}(s, a_i) = V^* + \epsilon_i, \quad \mathbb{E}[\epsilon_i] = 0, \quad \text{Var}(\epsilon_i) = \sigma^2
\end{equation}

\subsection{Standard Q-Learning Bias}
In standard Q-learning, the target value takes the maximum across action estimates:
\begin{equation}
    \mathbb{E}\left[ \max_i \hat{Q}(s, a_i) \right] = V^* + \mathbb{E}\left[ \max_{i=1,\dots,M} \epsilon_i \right]
\end{equation}
If $\epsilon_i \sim \mathcal{N}(0, \sigma^2)$, then by extreme value theory for independent Gaussians:
\begin{equation}
    \mathbb{E}\left[ \max_{i=1,\dots,M} \epsilon_i \right] \approx \sigma \sqrt{2 \ln M} > 0 \quad \forall M \ge 2
\end{equation}
This positive bias scales with $\sqrt{\ln M}$, causing standard DQN to systematically overestimate expected returns on candidate paths with noisy latency measurements.

\subsection{Double Q-Learning Decoupling}
In Double Q-learning, action selection is performed by set of weights $\theta_1$, while evaluation is performed by independent set $\theta_2$:
\begin{equation}
    a^* = \arg\max_i \hat{Q}(s, a_i; \theta_1)
\end{equation}
The evaluated target value is:
\begin{equation}
    \mathbb{E}\left[ \hat{Q}(s, a^*; \theta_2) \right] = \mathbb{E}\left[ V^* + \epsilon_2(a^*) \right] = V^* + \mathbb{E}[\epsilon_2(a^*)]
\end{equation}
Because $\theta_2$ is statistically independent of $\theta_1$, the estimation noise $\epsilon_2$ is independent of the index selection $a^*$:
\begin{equation}
    \mathbb{E}[\epsilon_2(a^*)] = \sum_{k=1}^M \mathbb{P}(a^* = k) \cdot \mathbb{E}[\epsilon_2(k)] = \sum_{k=1}^M \mathbb{P}(a^* = k) \cdot 0 = 0
\end{equation}
Thus, Double Q-learning completely eliminates maximization overestimation bias, proving theoretical stability.

\section{Derivation of Importance Sampling Correction in PER}
Under Prioritized Experience Replay, transitions are drawn from distribution $P(i) = p_i^\alpha / \sum p_k^\alpha$ instead of the true empirical replay distribution $U(i) = 1/N$.

The expected value of gradient updates under the true uniform distribution is:
\begin{equation}
    \mathbb{E}_{i \sim U}[\nabla_\theta \mathcal{L}_i] = \sum_{i=1}^N \frac{1}{N} \nabla_\theta \mathcal{L}_i
\end{equation}
Under prioritized sampling distribution $P(i)$, the expectation is:
\begin{equation}
    \mathbb{E}_{i \sim P}\left[ w_i \nabla_\theta \mathcal{L}_i \right] = \sum_{i=1}^N P(i) \cdot w_i \nabla_\theta \mathcal{L}_i
\end{equation}
Equating the two expectations to eliminate bias:
\begin{equation}
    P(i) \cdot w_i = \frac{1}{N} \implies w_i = \frac{1}{N \cdot P(i)}
\end{equation}
Raising to exponent $\beta$ and normalizing by $\max_j w_j$ yields the standard PER Importance Sampling weight formula:
\begin{equation}
    \tilde{w}_i = \left( \frac{1}{N \cdot P(i)} \right)^\beta \Big/ \max_j w_j
\end{equation}
When $\beta = 1.0$, $\mathbb{E}_{i \sim P}[\tilde{w}_i \nabla \mathcal{L}_i] \propto \mathbb{E}_{i \sim U}[\nabla \mathcal{L}_i]$, guaranteeing that prioritized stochastic gradient descent converges to the true expected value.

\section{Derivation of the M/M/1 Queuing Delay Model}
Consider an OpenFlow switch output port buffer modeled as a birth-death Markovian queuing system $M/M/1$:
\begin{itemize}
    \item Arrival rate: $\lambda$ packets/second (Poisson process).
    \item Service rate: $\mu = C / \bar{L}$ packets/second (exponentially distributed service times).
    \item Traffic intensity (link utilization): $\rho = \lambda / \mu < 1$.
\end{itemize}

Let $P_n$ denote the steady-state probability that exactly $n$ packets reside in the queue buffer. The balance equations state:
\begin{equation}
    \lambda P_{n-1} = \mu P_n \implies P_n = \left(\frac{\lambda}{\mu}\right)^n P_0 = \rho^n P_0
\end{equation}
Since $\sum_{n=0}^\infty P_n = 1$:
\begin{equation}
    P_0 \sum_{n=0}^\infty \rho^n = \frac{P_0}{1 - \rho} = 1 \implies P_0 = 1 - \rho
\end{equation}
Thus, $P_n = (1 - \rho)\rho^n$. The mean number of packets in the system $L$ is:
\begin{equation}
    L = \sum_{n=0}^\infty n P_n = (1 - \rho) \sum_{n=0}^\infty n \rho^n = (1 - \rho) \frac{\rho}{(1 - \rho)^2} = \frac{\rho}{1 - \rho}
\end{equation}
Applying Little's Law ($L = \lambda W$), the average time $W$ a packet spends in the system (transmission delay + queuing delay) is:
\begin{equation}
    W = \frac{L}{\lambda} = \frac{\rho}{\lambda(1 - \rho)} = \frac{1}{\mu(1 - \rho)} = \frac{1}{\mu - \lambda}
\end{equation}
Subtracting the base transmission delay $1/\mu$, the average queuing delay $W_q$ is:
\begin{equation}
    W_q = W - \frac{1}{\mu} = \frac{1}{\mu - \lambda} - \frac{1}{\mu} = \frac{\lambda}{\mu(\mu - \lambda)} = \frac{\rho}{\mu(1 - \rho)}
\end{equation}
This establishes the theoretical foundation for Equation~\ref{eq:queuing_delay_formula}.
